% ==============================================================================
% 第 30 章：总结 (Chapter 30: Summary)
% ==============================================================================

\chapter{总结}
\label{sec:ch30_summary}

本书第五部分将“顿悟”（Eureka!）的灵感引入了程序设计。与前四部分严格遵循数据定义的结构化设计不同，灵感设计始于对程序应\textbf{如何解决问题}或\textbf{如何处理表征问题的数据}的算法巧思。在这一语境下，设计意味着构想出一种精妙机制：在与原问题性质相同但规模更简单的新子问题上发起递归调用。

请牢记：尽管我们将其学术性地命名为\term{生成型递归（generative recursion）}，但绝大多数计算机科学家与软件工程师通常直接将此类函数称为\term{算法（algorithms）}。

学完本书的这一部分之后，你应该对生成型递归的设计建立起如下深刻理解：

\begin{enumerate}
  \item \textbf{设计配方的大体框架依然坚挺有效}：问题分析、数据定义、签名、目的说明、示例用例与测试步骤一如既往地规范着我们的思维路线。
  \item \textbf{主要变革发生于从模板到编码的实现步骤}：通用算法蓝图引出了关于如何具体实现函数的\textbf{四个核心问题}。其中前两个问题用于确立并解决“平凡”（trivial）边界情形，而后两个问题则用于规划新子问题的生成策略以及子问题之解的拼装合并。
  \item \textbf{关键的增补在于终止性分析}：与天然保证终止的结构型递归函数不同，算法在某些特定输入下可能无法停机。这种非终止现象可能源于算法思想本身的理论局限性，也可能源于从数学到代码实现过程中的离散化缺陷（如除零、浮点溢出或图环路死循环）。无论何种情况，未来的代码读者与维护者都理应在目的说明中获得关于潜在异常输入的清晰警示。
\end{enumerate}

在真实的职业编程生涯中，你会频繁遇到一些简单或经典的现成算法，并被期望能够轻车熟路地将其落地实现。而对于那些真正极其复杂精妙的尖端算法，顶尖科技公司通常会聘请高薪专家、领域学者和数学家在白板上推导完所有数学细节后，再交由程序员将其转化为高效软件系统。你也必须为此类高难度协作做好充分准备，而通往卓越的最佳途径唯有躬行实践。
